% !PRINTED 571
\weberpage{571}
%<<CONT
[\S.\,169, (6)] thereby takes on another form. We thus obtain
the same function $\psi_\lambda(\alpha)$ when we replace $\lambda,\mu$ by $\lambda f, f$ or
by $f, \lambda f$, hence [\S.\,169, (8)]

% [resolved 2026-08-06; confirmed 2026-08-07 at 600 dpi] CORRECTED THROUGHOUT THIS LEAF.
% All four sigmas in the two displays carry "t" ABOVE only; what sits beneath each is a short
% horizontal RULE under the sigma, not a second "t". The original transcription "\sum_t^t"
% duplicated the index. Reproduced here as "t" above, with the printed rule recorded in this
% note rather than faked in LaTeX.
\[
\psi_\lambda(\alpha) = \sum^{t} \alpha^{\operatorname{ind} t - (\lambda+1)\operatorname{ind}(t+1)}
= \sum^{t} \alpha^{\lambda \operatorname{ind} t - (\lambda+1)\operatorname{ind}(t+1)}.
\]

If, then, two numbers $\lambda,\lambda'$ are connected with one another by the
congruence

\weq{1}{
\lambda\lambda' \equiv 1 \pmod{e}
}

it follows that

\[
\psi_\lambda(\alpha^{\lambda'}) = \sum^{t} \alpha^{\lambda\lambda' \operatorname{ind} t - (\lambda\lambda'+\lambda')\operatorname{ind}(t+1)}
= \sum^{t} \alpha^{\operatorname{ind} t - (\lambda'+1)\operatorname{ind}(t+1)},
\]

hence the formula

\weq{2}{
\psi_\lambda(\alpha^{\lambda'}) = \psi_{\lambda'}(\alpha).
}

If we then put $-\lambda-\mu$ for $\lambda$ in \S.\,169, (5), and
hence $-\lambda$ for $\lambda+\mu$ (where it must then be assumed that
$\lambda$ is not divisible by $m$), it follows that

\[
\frac{(\varepsilon^{-\lambda-\mu}, r)(\varepsilon^\mu, r)}{(\varepsilon^{-\lambda}, r)} = \psi_{-\lambda-\mu,\mu}(\varepsilon).
\]

We multiply the left side by

\[
\frac{(\varepsilon^{\lambda+\mu}, r)(\varepsilon^\lambda, r)}{(\varepsilon^{\lambda+\mu}, r)(\varepsilon^\lambda, r)} = 1,
\]

and apply in the numerator and denominator the formula

\[
(\varepsilon^\lambda, r)(\varepsilon^{-\lambda}, r) = (-1)^\lambda n
\]

[\S.\,169, (3)]; then it follows that

\[
(-1)^\mu \frac{(\varepsilon^\lambda, r)(\varepsilon^\mu, r)}{(\varepsilon^{\lambda+\mu}, r)} = \psi_{-\lambda-\mu,\mu}(\varepsilon),
\]

or

\[
\psi_{-\lambda-\mu,\mu}(\varepsilon) = (-1)^\mu \psi_{\lambda,\mu}(\varepsilon),
\]

and if we replace $\lambda,\mu$ by $\lambda f, f$ and determine $\lambda''$ from the
congruence

\weq{3}{
\lambda + \lambda'' + 1 \equiv 0 \pmod{e}
}

we obtain

\weq{4}{
\psi_\lambda(\alpha) = (-1)^f \psi_{\lambda''}(\alpha).
}

This formula can then in turn be combined with formula (2),
and the number of functions $\psi_\lambda$ to be computed for a given
$e$, whose number amounts to $e-2$, thereby be reduced to about
the sixth part.

A third formula results when one puts $\varepsilon^{-1}$ for
$\varepsilon$ in \S.\,169, (5) and then multiplies, thus
%CONT>>
